\newcommand{\esameinfo}{
    \chapter*{Informazioni sull'esame}
    \addcontentsline{toc}{chapter}{Informazioni sull'esame}
    
    \section*{Modalità d'esame}
    
    L'esame è composto da due parti:
    \begin{itemize}
        \item \textbf{Scritto}
        \item \textbf{Laboratorio di Processing}
    \end{itemize}
    
    \subsection*{Scritto (da 18 a 31 punti)}
    Consiste in due fasi che si svolgono una dopo l'altra:
    
    \begin{enumerate}
        \item Test a risposta multipla con 10 domande. Si supera con 6/10. Questo numero di risposte corrette si chiama \textbf{A}.
        \item Chi supera la fase 1 svolge un test a risposta aperta di tre domande. L'esito darà un punteggio compreso tra 0 e 15. Chiamiamo \textbf{B} questo punteggio.
    \end{enumerate}
    
    Il voto finale dello scritto sarà uguale a:
    \[
    \text{Voto scritto} = 1.6 \times A + B
    \]
    
    \subsection*{Laboratorio di Processing}
    Verranno assegnati degli esercizi da svolgere in un tempo ragionevole, tramite i PC del laboratorio.
    
    \subsection*{Formula di valutazione del voto finale}
    \[
    \text{Voto finale} = \frac{(\text{voto scritto} \times 6) + (\text{voto lab} \times 3)}{9}
    \]
    
    \vspace{0.5cm}
    
    \begin{center}
        \large
        \textbf{Teams}: \texttt{1vd4wnf}
    \end{center}
}